\clearpage\thispagestyle{empty}
\pdfbookmark[0]{Quelle zum Brief an von Lindenau vom 25. Januar 1808}{lindenau1808-sources}
\begin{center}\textbf{Quelle zum Brief an von Lindenau vom 25. Januar 1808}\end{center}
\noindent Diese KI-Übertragung ins heutige Deutsch folgt dem Brief von Carl Friedrich Gauß an von Lindenau vom 25. Januar 1808 in \href{https://gdz.sub.uni-goettingen.de/id/PPN236007467}{Gauß’ \textit{Werken}, Bd.~VI} (Göttingen, 1874), S.~296--298. Der Brief erschien zuerst als \textit{Auszug aus einem Schreiben des Herrn Professors Gauss} in der \textit{Monatlichen Correspondenz}, Bd.~XVII, Februar 1808, \href{https://www.digitale-sammlungen.de/en/view/bsb10538608?page=188}{S.~182}--\href{https://www.digitale-sammlungen.de/en/view/bsb10538608?page=194}{188}. Das digitalisierte Exemplar stammt aus der Bayerischen Staatsbibliothek.
\par\medskip
Die Werkausgabe verweist auf die Kometenbeobachtungen in der Anzeige vom 25. Februar und stellt die Mailänder Vestabeobachtungen und ihre Differenzen nebeneinander. Im Erstdruck stehen sie auf \href{https://www.digitale-sammlungen.de/en/view/bsb10538608?page=189}{S.~183} und \href{https://www.digitale-sammlungen.de/en/view/bsb10538608?page=190}{S.~184}. In der Pallastafel bleibt das hochgestellte Zeichen nach 770 aus der Werkausgabe erhalten. Im Erstdruck stehen dort ein Komma und ein Sekundenzeichen. Der Brief endet in der Werkausgabe mit einer in eckige Klammern gesetzten Beschreibung der Juno-Ephemeride. Die Ephemeride steht in der Zeitschrift auf \href{https://www.digitale-sammlungen.de/en/view/bsb10538608?page=192}{S.~186}--\href{https://www.digitale-sammlungen.de/en/view/bsb10538608?page=194}{188}. KI-Übertragung und KI-Satz.
\clearpage
